\section{The model}
\label{sec:model}

\subsection{Agents, states, actions, and consequences}
\label{sec:primitives}

There are $n$ agents, $N=\{1,\dots,n\}$.  Time is divided into generations
$g=1,2,\dots$, each consisting of $T$ episodes.  In each episode a state
$s\in S$, with $S$ finite, is drawn independently from a full-support
distribution $\rho\in\Delta(S)$; each agent takes an action $a_i$ from a
finite set $A_i$ through a committed policy described below; and a
consequence profile $x=(x_1,\dots,x_n)$ is realized in a finite space
$X=\prod_iX_i$.  Agent $i$'s consequence $x_i$ is everything the episode
delivers to it---profit, sales, observed reactions---and its payoff is a
function $u_i:X_i\to\mathbb R$ of that consequence.  Write $A=\prod_iA_i$ for
action profiles.

States are the analyst's descriptions of the payoff-relevant situations
agents may face: customer configurations, input conditions, competitive
circumstances.  Their status in the agents' own representations is the
subject of \cref{sec:awareness-sub}.

\subsection{Mechanisms and the objective causal system}
\label{sec:mechanisms-sub}

\begin{definition}[Elementary mechanisms and the objective system]
\label{def:mechanisms}
A finite set $\Theta$ of \emph{elementary causal mechanisms} is given, each a
map $\theta:A\to\Delta(X)$.  The \emph{objective assignment} is a function
$\tau:S\to\Theta$, and the \emph{objective causal system} is
\begin{equation}
F(\cdot\mid s,a)=\tau(s)(a)\in\Delta(X).
\label{eq:objective-system}
\end{equation}
Two states are \emph{causally equivalent} if they carry the same mechanism;
$\mathcal P^{\tau}$ denotes the partition of $S$ into causal-equivalence
classes, and $\marg{i}\theta(a)$ the $X_i$-marginal of $\theta(a)$.
\end{definition}

The causal reading of \cref{eq:objective-system} is maintained throughout and
does real work.  $F$ is a response function, not a statistical summary: it
specifies the distribution of consequences that \emph{would} result from
every action profile at every state, including profiles never played.
Agents' actions, together with the state, cause their consequences.  A
mechanism is the reduced form of whatever detailed model---a structural
causal model of customer behavior, a market microstructure, a production
process---underlies that relation; the paper imposes no structural-causal-model
formalism, but examples may exhibit micro-foundations explicitly, as the
running example's customer market does.  Three consequences of the causal
reading recur.  First, theories, defined next, are interventional objects:
they make predictions at unplayed profiles.  Second, on-path data constrain
$F$ only at visited state--profile pairs, so observationally equivalent
systems with different counterfactual content are generic; that is the engine
of the identification results.  Third, the on-path distribution of
consequences given an agent's category of states is a \emph{selection}
object---it mixes mechanisms according to which states co-occur with which
play---and must never be conflated with any mechanism's own conditional.

\subsection{Awareness, theories, and conjectures}
\label{sec:awareness-sub}

\begin{definition}[Awareness and expressible theories]
\label{def:awareness}
Agent $i$'s awareness partition is $\mathcal P_i\in\Pi(S)$, the set of
partitions of $S$.  Every consciously available object---theory, belief,
conjecture, policy, question---must be measurable with respect to
$\mathcal P_i$.  The \emph{expressible theories} are the cell-constant
mechanism assignments
\begin{equation}
M(\mathcal P_i)=\{m:\mathcal P_i\to\Theta\},
\label{eq:expressible}
\end{equation}
with $m$ predicting $x_i\sim\marg{i}m(C)(a)$ at cell $C$ under profile $a$.
Partition $\mathcal P$ is \emph{sufficient} for an assignment $\tau'$ if
$\tau'$ is constant on every cell of $\mathcal P$, equivalently if $\mathcal
P$ refines $\mathcal P^{\tau'}$; the objective assignment is expressible for
agent $i$ if and only if $\mathcal P_i$ is sufficient for $\tau$.
\end{definition}

An awareness cell is thus an implicit causal claim: the states it pools are
alike in how actions produce consequences.  An expressible theory makes the
claim explicit by naming the mechanism.  Restricting theories to
\emph{single} mechanism assignments---rather than arbitrary conditional
distributions per cell---is substantive and deliberate.  If any distribution
were expressible, any within-cell mixture would be too, and no experience
could reject the class: coarseness alone cannot produce a detectable failure
in a saturated model space.  The mechanism vocabulary $\Theta$ is known to
all agents from the outset; what an agent may lack is the \emph{distinction}
that separates states carrying different mechanisms.  Everything the paper
calls insight is the production of such a distinction.  Expansion of the
mechanism vocabulary itself is a deeper dimension of awareness growth
deliberately left outside this paper.

Agent $i$ holds a belief $\mu_i\in\Delta(M(\mathcal P_i))$ over expressible
theories and a conjecture $q_i:\mathcal P_i\to\Delta(A_{-i})$ about rival
play, both $\mathcal P_i$-measurable.  Theories and conjectures are different
kinds of objects and are deliberately separated: a theory is a causal claim
about the objective system; a conjecture is a strategic claim about what
other agents will do.  The diagnostic apparatus developed below tests
theories, not conjectures.

\begin{assumption}[Observation and retention]
\label{ass:observation}
At the end of each episode, agent $i$ observes
$(C_{\mathcal P_i}(s),a,x_i)$: its awareness cell containing the realized
state, the realized action profile, and its own consequence.  Episodes are
retained as tokens in a deliberative record $r_i$ and can be retrospectively
recoded after the agent acquires a new distinction; before acquisition,
neither belief nor action can condition on that distinction.  The
analyst-level history additionally records $s$ itself, so that a newly
produced representation can recode the tokens.
\end{assumption}

Retention without recodability-in-advance is the minimal separation between
experiencing a particular and possessing a projectible concept for it: the
agent can remember ``customers like that one'' without commanding a predicate
that distinguishes them, and the technology of \cref{sec:inquiry} operates on
exactly this material.  Observability of the full action profile keeps the
separation of causal from strategic error clean; a variant folding components
of $a$ into $x_i$ changes bookkeeping only.

\begin{assumption}[Commitment and subjective optimality]
\label{ass:commitment}
At the start of each generation, agent $i$ commits to a policy
$\sigma_i:\mathcal P_i\to A_i$ satisfying
\begin{equation}
\sigma_i(C)\in\arg\max_{a_i\in A_i}
\sum_{m\in M(\mathcal P_i)}\mu_i(m)
\sum_{a_{-i}\in A_{-i}}q_i(a_{-i}\mid C)\,
\mathbb E_{x\sim m(C)(a_i,a_{-i})}\!\left[u_i(x_i)\right],
\label{eq:subjective-optimality}
\end{equation}
holds it for the generation's $T$ episodes, and revises beliefs, conjectures,
and policies only between generations.  When dissonance
(\cref{sec:dissonance}) has made $\mu_i$ undefined and no new representation
has been produced, the previously committed policy is carried forward.
Agents evaluate one generation ahead; no discounting is imposed.
\end{assumption}

Commitment is where statistics and economics meet: a test needs a stable
data-generating configuration, and a generation of committed play supplies
exactly that.  What a deterministic model must assume as a batching device
is, in the stochastic model, the natural granularity of learning.

\subsection{On-path conditionals}

Fix a generation with committed profile $\sigma=(\sigma_1,\dots,\sigma_n)$.
Each state $s$ receives the deterministic profile
$a(s)=(\sigma_j(C_{\mathcal P_j}(s)))_{j\in N}$.  Agent $i$'s \emph{on-path
conditional} at an observation pair $(C,a)$ with
$S(C,a)=\{s\in C:a(s)=a\}\neq\emptysetset$ is
\begin{equation}
p_i^{*}(\cdot\mid C,a)
=
\sum_{s\in S(C,a)}\rho\bigl(s\mid S(C,a)\bigr)\,\marg{i}\tau(s)(a).
\label{eq:onpath-conditional}
\end{equation}
This is the selection object flagged above: rivals with finer partitions
condition on distinctions agent $i$ cannot see, so the within-cell state
mixture can vary with the profile.  It is precisely this
dependence---differentiated play by others inside one's own
categories---that generates the contrasts on which the diagnostic apparatus
feeds.

\subsection{Modeling commitments}
\label{sec:commitments}

Five disciplines govern everything that follows; they are what distinguish a
model of unawareness from a model of inattention or model selection.
(i)~Agents hold no prior over the outputs of the inquiry technology and never
optimize over named future theories; pre-insight deliberation quantifies only
over represented objects.  (ii)~Questions are generated by defects in the
current representation, not drawn from a primitive list.  (iii)~Inquiry can
fail, and its successes can be false; nothing reveals the objective
assignment by construction.  (iv)~All valuations attached to inquiry are
subjective, coarse, and possibly wrong; no objective payoff information
leaks from behind the awareness frontier.  (v)~Evidence that generates a
candidate representation is separated, in the formal accounting, from
evidence that warrants it.

\subsection{The running example}

\begin{example}[A four-state laboratory]
\label{ex:setup}
Two firms serve a unit market.  States are customer configurations
$S=\{s_{00},s_{01},s_{10},s_{11}\}$, displayed as a $2\times2$ array with row
coordinate $r$ and column coordinate $c$; $\rho$ is uniform.  Actions are
product architectures $A_i=\{0,1\}$.  The mechanism vocabulary is
$\Theta=\{\theta^0,\theta^1,\theta^n\}$, with a demand parameter
$\beta\in[0,1]$: under $\theta^a$ (``architecture $a$ is preferred''), if the
firms differentiate, the firm offering $a$ receives payoff $1$ and its rival
$0$; if they pool on $a$, each receives $\tfrac12$; if they pool on the other
architecture, each receives $\beta/2$, the customer withholding $1-\beta$ of
the unit value.  Under $\theta^n$ (``architecture is neutral''), each firm
receives $\tfrac12$ regardless.  A consequence $x_i$ is the pair (own payoff,
observed action profile).  The objective assignment is the \emph{row law}
$\tau(s_{rc})=\theta^{r}$: the row coordinate is causally relevant, the
column coordinate is not.  Both firms begin with the coarse partition
$\mathcal P_i^0=\{S\}$, under which the expressible theories are the three
constant assignments; the row law is not among them, because it is not
constant on $S$.  The baseline sets $\beta=1$ (deterministic winner-take-all
with market splitting); $\beta<1$ returns in \cref{sec:dynamics}.  All
mechanisms here are deterministic, so every test below operates in its exact,
zero-tolerance instance; nothing in the general development relies on that.
\end{example}
